% supp_protocol2.tex -- Supplementary Section S10: protocols of the second outbreak test and of the further analyses; register of all tests.
% Paths in "% src:" comments are relative to the repository root.
% Numbers: data/v2_numbers.json and data/v2_tab_*.tex (code/v2_numbers.py).
\section{Protocols of the second outbreak test and of the further analyses; register of all tests}\label{appS_prot2:sec}

This section documents the four analyses of Section~\ref{M-s8_outbreaks:sec} of the main text other than the first
outbreak test: what was fixed in advance and the evidence for when; the decision rules and their
power; the data and what may be redistributed; and a register of every test of Section~\ref{M-s8_outbreaks:sec}, which separates what
was pre-specified from what was added afterwards, by the analysis itself or by its re-check. The four
protocols, their power addenda, freeze records and post-hoc logs are released with the code, under the file
names they were given (\texttt{PREREGISTRATION.md} and so on), as redacted copies (Section~\ref{appC_protocol:sec});
as in
Section~\ref{appC_protocol:sec}, they are called protocols here and their contents pre-specified. `The
re-check' of an analysis is its re-implementation with separately written code, within the same project and
on the same machine, as described in Section~\ref{M-s1_intro:sec}; where it corrected a result, the corrected result is the one reported.
Where the frozen protocol texts and logs speak of a `verifier', they mean this
same-project re-check, not an independent party, and by `the analyst' they mean the first analysis (both
passes were made by the author). The
frozen protocol texts call the four analyses avenues 1 to 4, and the first outbreak test with its dataset and the
later analyses rounds 1 and 2 (two successive sets of analyses); the article does not use these labels.

\subsection{What was frozen, and the evidence for when}\label{appS_prot2:sec:freeze}

\begin{table}[tbp]
  \centering\footnotesize
  \caption{The four protocols of the second outbreak test and of the further analyses. Times are British Summer Time on 1~October 2026, from the
  freeze records (self-reported) and from the modification times of the files. Hashes: number of files whose
  SHA-256 is stored in the freeze records; all were recomputed for this article on the unredacted originals
  and match.}
  \label{appS_prot2:tab:freeze}
  % src: data/v2_numbers.json (freeze) <- code/bsc_validation2/a*/PREREG_FREEZE.json, FIX_FREEZE.json, PREREG_FREEZE_2_before_unblinding.json, PREREG_EXT_FREEZE*.json; file times of the working tree
  \setlength{\tabcolsep}{4pt}
  \begin{tabular}{@{}>{\raggedright\arraybackslash}p{0.12\textwidth}>{\raggedright\arraybackslash}p{0.21\textwidth}c>{\raggedright\arraybackslash}p{0.24\textwidth}>{\raggedright\arraybackslash}p{0.24\textwidth}@{}}
    \toprule
    Analysis & Frozen & Hashes & First result that uses outcomes & Known beforehand \\
    \midrule
    Further outbreaks (Section~\ref{M-s8_outbreaks:sec:second}) &
      protocol 09:42:01; calibration and power 09:52:41 & 4 + 6 &
      held-out strata scored 09:53:08 &
      every count, including the held-out ones; the outcome of the first test \\
    Contact records (Section~\ref{M-s8b_evidence:sec:contacts}) &
      protocol 09:23:27; modified walks 10:34:24, before the confirmatory datasets & 8 + 26 &
      first dataset 09:24:00 (the run takes 19\,s); first confirmatory dataset 10:35:34 &
      raw counts per day and group labels; the published properties of such data; expectation of failure
      written down \\
    Tracer measurements (Section~\ref{M-s8b_evidence:sec:tracer}) &
      protocol 09:18:16; kernels, predictions and power 09:44:43 & 23 + 10 &
      evaluation 09:49:08 &
      the outbreak counts and the failures of the first test; the tracer sources had been opened \\
    \quad extension to four flights &
      protocol 10:24:56; predictions 10:26:01 & 7 + 4 &
      evaluation 10:26:09 &
      the near/far counts and the result of the second outbreak test: outcome-aware \\
    Zone level (Section~\ref{M-s8b_evidence:sec:zones}) &
      protocol 09:48:58 & 14 &
      schools 09:56:44 (run 09:49--09:56); floors 10:21:08 &
      total of cases; qualitative conclusion of the source; for the floors, the counts \\
    \bottomrule
  \end{tabular}
\end{table}

Table~\ref{appS_prot2:tab:freeze} gives the times. The limits are those of the first test
(Section~\ref{appC_protocol:sec:freeze}) and apply to all four analyses.
(i)~Every protocol is \emph{data-aware} and \emph{model-blind}. What was fixed in advance is the list of
events or datasets, the models, the split into calibration and hold-out, the metrics, the decision rule and
the list of sensitivity analyses, each with a power analysis; the data themselves had been read.
(ii)~The freezes are self-attested. The records are files on the author's machine, with matching hashes and
no external time stamp or version-control history; the hashes show that the files have not changed since, not
when they were written. In the outbreak analysis the record was itself rewritten when the power addendum was
added, and the clock times written into its post-hoc log are one to two hours ahead of the file times.
(iii)~Parts of the code were written after the freeze: in the outbreak analysis the modules that build the
events and run the comparison (the last one four seconds before the hold-out output), in the tracer analysis
the evaluation script. The intervals are short because each analysis was prepared and run in one or two continuous stretches of
work.
(iv)~In the outbreak analysis the constant-ratio model was coded as $p_j=1-\exp(-\varepsilon\varrho)$ in the near
zone, not as the $p_j=\varrho\,p_0$ of the protocol. With the form of the protocol the pooled
$\Delta_{\mathrm{C}}$ is $-7.19$ in the re-check's re-implementation, against $-9.20$; the outcome is the
same.
(v)~The extension of the tracer analysis to four flights was written, run and evaluated within 73 seconds,
with the outcomes known; its freeze is nominal, and it is reported as exploratory throughout.
% src: notes/VERIFICATION.md (Section 4.6, second outbreak test, item 1; Section 4.7, contact records, item 1; Section 4.8, tracer measurements, items 1, 13; Section 4.9, zone level, item 1); data/v2_numbers.json (out2.check.registered_form.all: deltaC -7.194; freeze: extension protocol 10:24:56, evaluation 10:26:09)

\subsection{Further outbreaks: rule, power, numerical correction}\label{appS_prot2:sec:outbreaks}

\begin{table}[tbp]
  \centering\footnotesize
  \caption{(a)~Decision rule of the second outbreak test, applied to $\Mtwo$ (and, in separate columns of the
  analysis, to $\Mone$ in closed form and to an exploratory variant with a well-mixed fraction). (b)~Probability
  of each criterion and of the full claim V1 when the data are generated by the model in the first column, at
  its calibrated posterior (200{,}000 replicates), as computed before any held-out stratum count was scored
  and, in brackets, after the numerical correction of the room kernel.}
  \label{appS_prot2:tab:rule}
  % src: code/bsc_validation2/a1_more_outbreaks/PREREGISTRATION.md section 7; code/bsc_validation2/a1_more_outbreaks/src/a1/pipeline.py (V5 if two or more failures; elif not B1: V4; elif mirror: V3; elif B2 and no failure: V1; else V2); data/v2_numbers.json (out2.power.registered, out2.power.corrected (M2: B1B2 0.926 - V1 0.762 = 0.164), out2.power.registered_air)
  \begin{tabular}{@{}l>{\raggedright\arraybackslash}p{0.40\textwidth}>{\raggedright\arraybackslash}p{0.46\textwidth}@{}}
    \toprule
    \multicolumn{3}{@{}l}{(a) \emph{Criteria and outcomes}} \\
    B1 & better than well mixed & $\Delta_0>0$ with one-sided $p<0.05$ under $\Mzero$ \\
    B2 & better than the constant ratio & $\Delta_{\mathrm{C}}>0$ with one-sided $p<0.05$ under the
         constant-ratio model; mirror: $\Delta_{\mathrm{C}}<0$ with $p<0.05$ under $\Mtwo$ \\
    B3 & adequacy & two-sided predictive $p\ge0.05$ for every held-out event \\
    \midrule
    V1 & B1, B2, B3, and power of (B1 and B2) at least 0.8 &
         `predicts held-out distance gradients better than both a well-mixed model and a constant near/far
         risk ratio' \\
    V2 & B1 and B3 with at most one failure; neither B2 nor its mirror &
         `better than well mixed; not distinguishable from a constant ratio' \\
    V3 & B1; mirror of B2 & `better than well mixed, but a constant ratio predicts better' \\
    V4 & B1 not met & `no out-of-sample support for the kernel' \\
    V5 & two or more failures of B3 & a list of the events reproduced and of those not \\
    \multicolumn{3}{@{}p{0.95\textwidth}@{}}{\emph{Not covered:} B1 and B2 met with exactly one B3 failure was given no
    outcome by the protocol (probability 0.16 if $\Mtwo$ is true); the code classes it with V2. The observed
    outcome, V5, does not depend on this.} \\
    \bottomrule
  \end{tabular}

  \medskip
  \begin{tabular}{@{}lcccc@{}}
    \toprule
    (b) Truth & B1 & B2 & B1 and B2 & V1 \\
    \midrule
    well mixed       & 0.007 (0.015) & 0.096 (0.169) & 0.005 (0.010) & 0.003 (0.008) \\
    constant ratio   & 0.652 (0.732) & 0.029 (0.052) & 0.028 (0.049) & 0.020 (0.039) \\
    $\Mtwo$          & 0.989 (0.982) & 0.961 (0.935) & 0.956 (0.926) & 0.784 (0.762) \\
    \midrule
    $\Mtwo$, four flights only & 0.980 & 0.822 & 0.821 & -- \\
    \bottomrule
  \end{tabular}
\end{table}

The rule of Table~\ref{appS_prot2:tab:rule} was written to remove the weakness found in the first test: a
constant-ratio world passes B1 with probability 0.65 to 0.73, so that B1 alone is no evidence for the
kernel, and passes B1 and B2 together with probability 0.03 to 0.05. The re-check recomputed the
power from 400{,}000 draws of its own (0.927, 0.048 and 0.010 for B1 and B2 together under the three truths;
0.822 for the four flights).
% src: data/v2_numbers.json (out2.power.check: all M2 P_B1B2 0.9268, CRR 0.0476, M0 0.0100; air M2 0.8222)

\paragraph{The numerical correction.}\label{appS_prot2:sec:numerics}
The propagator of the first test writes $P_t(y|x)$ on a rectangle with reflecting walls as a sum over cosine
modes. Its absolute error reaches about $10^{-10}$ of the largest exposure in an event, so that receivers
whose true exposure is smaller than that carry round-off noise of either sign. Replacements by the corrected
propagator occur for mixing coefficients up to about 2.5 to $16\sqmh$, in the cabins and in the rooms of the
second test alike. The posterior of
the room coefficient first computed (`mode $0.02\sqmh$, bimodal') came from two such noise values in the
likelihood of the ward inpatients; it was found by inspection of a diagnostic plot after the held-out events
had been scored. In the corrected code each one-dimensional propagator is a positive sum of image terms, and
the image value replaces the mode sum wherever the latter is below $10^{-8}$ of the event maximum; where both
are resolved, at coefficients at which a replacement occurs, they agree to about $10^{-4}$ or better
(largest relative difference $8.4\times10^{-5}$, in the ward inpatients). The primary analysis, the power and the secondary analyses were then
re-run unchanged (the analyses stored in \texttt{10\_corrected.json}); the 15 sensitivity variants were not,
and are reported with the kernel as first frozen. For the flights no reported
digit changes (their posterior lies above the affected range). For the rooms the posterior becomes $794\sqmh$
(302--4{,}390), the score of the processing line $-15.24$ instead of $-16.22$ and its predictive $p$
$1.4\times10^{-5}$ instead of $5\times10^{-6}$; the pooled scores are $+14.35$ and $-9.15$ instead of $+13.37$
and $-10.13$; and the outcome is V5 in both cases. Both versions are stored. The re-check confirmed
the defect and the correction with a third propagator, built by uniformisation and sharing no code with
either: it agrees with the image sum to $10^{-14}$ at small coefficients, while the mode sum returns
non-positive exposures at up to 32 of the 74 ward receivers. On the restaurant calibration of the first test
the same defect changes the log-likelihood by about $10^{-7}$ at most at every coefficient in the posterior range,
so that calibration is not affected. The hold-out predictions of the first test were affected in the
classroom, where the mode sum returns non-positive exposures at up to 11 of the 24 desks for coefficients of
$0.1\sqmh$ and below; they were re-run with the corrected propagator for this article
(Section~\ref{M-s8_outbreaks:sec:holdout}).
% src: data/bsc_validation2/a1_more_outbreaks/09_numerics_check.json (per event: last coefficient with frac_replaced > 0: 2.5 (F2, F3, F4, F6, P1), 6.3 (F1, W2), 15.8 (F5, F7, F8, W1); max_rel_diff_where_resolved at those coefficients up to 8.4e-5 (W2)); data/s8_first_test_exact.json (kernel.rows: spectral_min_rel down to -1.28e-10); code/checks/outbreaks2/t6_first_set_floor.json (restaurant: 1.1e-7)
% src: code/bsc_validation2/a1_more_outbreaks/POSTHOC_LOG.md (P5); data/v2_numbers.json (out2.pooled_registered.M2.all, out2.pooled_registered.P1: M2 logp -16.22, p 4.9e-6; out2.hold.P1.M2: logp -15.24, p 1.42e-5); notes/VERIFICATION.md (Section 4.6, second outbreak test, item 13: confirmed); data/s8_first_test_exact.json (kernel.rows: spectral_nonpositive up to 11 of 24)

\begin{table}[tbp]
  \centering\small
  \caption{Mixing coefficient preferred by each event of the second dataset taken alone (profile likelihood of the
  stratum counts given the total; interval where the log-likelihood is within 1.92 of its maximum; grid from
  $10^{-2}$ to $10^{4}\sqmh$), and gain in log-likelihood over the largest coefficient of the grid
  ($10^{4}\sqmh$), which in a cabin is close to, but not yet, the well-mixed model.}
  \label{appS_prot2:tab:Dprofiles}
  % generated: data/v2_tab_Dprofiles.tex (code/v2_numbers.py <- data/bsc_validation2/a1_more_outbreaks/10_corrected.json, preferred_D_corrected)
  \begin{tabular}{@{}llccc@{}}
    \toprule
    Event & Pathogen & $\Dair$ ($\mathrm{m^2\,h^{-1}}$) & Interval & Gain \\
    \midrule
    \tabinput{data/v2_tab_Dprofiles}
    \bottomrule
  \end{tabular}
\end{table}

\begin{figure}[tbp]
  \centering
  \includegraphics[width=\textwidth]{v2_D_profiles_en}
  % generated by code/v2_figures.py (D_profiles) from data/bsc_validation2/a1_more_outbreaks/10_corrected.json (preferred_D_corrected.*.curves)
  \caption{Log-likelihood of the mixing coefficient of $\Mtwo$ for each event of the second dataset taken alone,
  relative to the largest coefficient of the grid ($10^{4}\sqmh$). Solid: held-out events; dashed: calibration events. The thick line in (a)
  is the sum over the eight flights (common coefficient).}
  \label{appS_prot2:fig:Dprofiles}
\end{figure}

\subsection{Contact records: data and statistics}\label{appS_prot2:sec:contacts}

\begin{table}[tbp]
  \centering\footnotesize
  \setlength{\tabcolsep}{3.6pt}
  \caption{The six contact datasets (SocioPatterns, RFID proximity at 20\,s resolution) and four statistics of
  the test days, observed and under the homogeneous lattice walk calibrated on the training days. Role:
  development (dev.) or confirmatory (conf.) dataset for the modified walks. Days: training + test. Sites: size of the lattice fixed by the calibration. Exponent: of the number of contacts per
  20 minutes against the number of people present (not comparable for the primary school, where occupancy
  hardly varies; the conference has no groups).}
  \label{appS_prot2:tab:datasets}
  % generated: data/v2_tab_datasets.tex (code/v2_numbers.py <- data/bsc_validation2/a2_contact_mobility/{s1,s3a}_<dataset>.json)
  \begin{tabular}{@{}>{\raggedright\arraybackslash}p{0.19\textwidth}lccccccc@{}}
    \toprule
     & & & & & \multicolumn{4}{c}{observed / walk} \\
    \cmidrule(l){6-9}
    Setting & Role & People & Days & Sites & \shortstack{Partners\\per day} & \shortstack{Within-group\\share} &
      Burstiness & Exponent \\
    \midrule
    \tabinput{data/v2_tab_datasets}
    \bottomrule
  \end{tabular}
\end{table}

\begin{figure}[tbp]
  \centering
  \includegraphics[width=\textwidth]{v2_contacts_structure_en}
  % generated by code/v2_figures.py (contacts_structure) from data/bsc_validation2/a2_contact_mobility/*.json (obs, models.*.stats)
  \caption{Contact statistics of the test days: value under each model divided by the observed value (shaded:
  tolerance 0.8 to 1.25; values below 0.02 are drawn at the left edge). The lattice walk matches what it was
  calibrated on, the total contact time and the mean duration, and misses long contacts, the number of
  partners, the differences between people, the repetition of pairs and the group structure. The anchored
  variants were developed on the first three datasets (top row) after that failure.}
  \label{appS_prot2:fig:structure}
\end{figure}

Table~\ref{appS_prot2:tab:datasets} lists the datasets. The first three were used to develop the modified
walks and the last three were not touched by any model or statistic until the modifications and their fitting
code had been frozen. Clock offsets, the removal of days with fewer than 1\pct{} of the records and the
definition of a person's presence (from the first to the last own contact of the day) are fixed in the frozen
loader. The seven families of statistics, with their tolerances, are: total contact time (ratio in 0.8--1.25);
mean contact duration and share of contact time in contacts of at least five minutes; burstiness of the gaps
between a person's contacts (absolute difference at most 0.10) and recurrence of pairs; mean and coefficient
of variation of the number of partners per person and day; coefficient of variation of individual contact
time; share of a day's pairs already in contact on the previous day; share of contact time within the recorded
groups. A family is passed if all its statistics are.
% src: code/bsc_validation2/a2_contact_mobility/PREREGISTRATION.md sections 2, 5; code/bsc_validation2/a2_contact_mobility/FIX_FREEZE.json
The power analysis was made on synthetic data before the freeze, with four replicates per combination of
generating model and candidate: a true or nesting model passed every family and the epidemic criterion in 20
of 20 replicates, and a wrong non-nested model reached the per-dataset criterion in none of 44. The
re-check qualified this in three ways: the 44 cases are not independent; the analysis covers the
base models only, and the modified walks fail on their own synthetic truth
(Section~\ref{M-s8b_evidence:sec:contacts}); and on real data the tolerance of the epidemic criterion is at the
level of the noise, since two runs of 4{,}000 epidemics on the same recorded network differ by up to 7\pct{} in
the secondary cases of the index case, against a tolerance of 10\pct. The rule is therefore adequate for the
failure of the lattice walk, whose errors are 25 to 72\pct{}, and weak for anything close to the tolerance.
% src: code/bsc_validation2/a2_contact_mobility/PREREGISTRATION.md section 7; code/bsc_validation2/a2_contact_mobility/POSTHOC_LOG.md item 1; data/checks/contacts.json: power (20/20, 0/44), noise_of_epidemic_criterion (4,000 runs, up to 7 %, tolerance 10 %), lattice_walk_errors_pct (25, 72); notes/VERIFICATION.md (Section 4.7, contact records, item 9: partially confirmed)

\subsection{Tracer measurements: rule and power}\label{appS_prot2:sec:tracer}

\begin{table}[tbp]
  \centering\footnotesize
  \caption{Tracer analysis. (a)~Decision rule over seven events (coach, bus, flight, trains, restaurant,
  classroom, call centre); D1: number of events rejected at the 5\,\% level; D2: pooled log score above that of
  $\Mzero$ with $p<0.05$; D3: above that of the constant ratios 2 and 3, each with $p<0.05$. (b)~Probability of
  each outcome when the data are generated by the model in the first column, computed after the kernels had
  been estimated and before any comparison with outbreak counts.}
  \label{appS_prot2:tab:tracerrule}
  % src: code/bsc_validation2/a3_tracer_physics/PREREGISTRATION.md section 6; data/v2_numbers.json (tr.power.truths.*.outcomes)
  \begin{tabular}{@{}l>{\raggedright\arraybackslash}p{0.82\textwidth}@{}}
    \toprule
    S1 & no D1 failure, D2 and D3: `with the kernel fixed by physical measurement the framework predicts the
         spatial patterns, and better than a generic gradient' \\
    S2 & no D1 failure, D2, not D3: `compatible with all events and better than well mixed; not shown to be
         better than a generic gradient' \\
    S3 & no D1 failure, not D2: `compatible, not discriminating' \\
    S4 & exactly one D1 failure: a sentence that names it \\
    S5 & two or more D1 failures: the hypothesis fails \\
    \bottomrule
  \end{tabular}

  \medskip
  \begin{tabular}{@{}lccccc@{}}
    \toprule
    (b) Truth & S1 & S2 & S3 & S4 & S5 \\
    \midrule
    tracer kernel     & 0.213 & 0.117 & 0.453 & 0.196 & 0.021 \\
    well mixed        & 0.009 & 0.004 & 0.488 & 0.401 & 0.099 \\
    constant ratio 2  & 0.017 & 0.204 & 0.112 & 0.455 & 0.211 \\
    constant ratio 3  & 0.000 & 0.022 & 0.002 & 0.204 & 0.773 \\
    \bottomrule
  \end{tabular}
\end{table}

The rule (Table~\ref{appS_prot2:tab:tracerrule}) is safe against a false claim and weak at confirming a true
one: if the tracer kernel were true it would earn the full claim S1 with probability 0.21 only, because on the
coach and the bus it predicts almost what the well-mixed model predicts. The outcome obtained, S4, has
probability 0.20 under the kernel, 0.40 if outbreaks are well mixed and 0.46 or 0.20 under a constant ratio of
2 or 3, so that it does not discriminate. Every model fails on the train table, and the rule counts this as a
failure of the kernel without asking whether a competitor passes. For the restaurant the rule fails the kernel
only if it is rejected at every count from two to five cases, which is lenient; the main text reports the
result at each count.
% src: data/v2_numbers.json (tr.power.truths: P S1 0.213, S4 0.196; M0 S4 0.401; CRR2 S4 0.455; CRR3 S4 0.204); code/bsc_validation2/a3_tracer_physics/POSTHOC_LOG.md (P7); notes/VERIFICATION.md (Section 4.8, tracer measurements, item 14)
Two departures from the protocol were forced by the data and are logged: the forward cabin has unmasked
breathing releases at three seats, not four, so that the flight uses two measured kernels in place of three;
and the minibus of the Hunan outbreak has 15 passengers with a simulated tracer value, not the 17 of the first
test. A full re-run of the stored analysis into a separate directory
reproduced every result file without a difference; the re-check then re-derived the coefficients
with its own solver (identical grid optima) and re-typed the tracer values from the sources.
% src: code/bsc_validation2/a3_tracer_physics/POSTHOC_LOG.md (P1, P2, P8); notes/VERIFICATION.md (Section 4.8, tracer measurements, items 2, 3: confirmed)

\subsection{Zone level: models and power}\label{appS_prot2:sec:zones}

A susceptible pupil of class $c$ has onset on day $t$ with probability $1-\exp(-\lambda_c(t))$,
$\lambda_c(t)=\beta\,x_c(t)+\alpha_0+\alpha_1P_{\mathrm{city}}(t)$, where $P_A(t)$ is the kernel-weighted
number of earlier onsets in the set $A$ per pupil of $A$ and the last term is the same quantity for the other
schools of the city. The kernel is a gamma density with mean 1.7 and standard deviation 1.0 days on daily bins,
as in \citet{endo2021within}, who chose it to give the mean serial interval of 2.2 days of influenza A(H3N2)
\citep{vink2014serial}. The three sets are exclusive: for a pupil of class $c$, $P_{\mathrm{class}}$ counts the
other pupils of the class (denominator $n_c-1$), $P_{\mathrm{grade}}$ the pupils of the other classes of the
grade and $P_{\mathrm{school}}$ the pupils of the other grades. The zone model has
$x_c=0.762\,P_{\mathrm{class}}+0.104\,P_{\mathrm{grade}}+0.134\,P_{\mathrm{school}}$, with the shares measured
in Lyon; the well-mixed school has $x_c=P_{\mathrm{whole\ school}}$; independent classes have
$x_c=P_{\mathrm{class}}$; the generic model has one fitted ratio $\varrho$ of own-class to rest-of-school risk;
and the ceiling has three free shares. Within-school terms are switched off during the winter break.
In the 36 grades that have a single class (453 pupils) the grade term is set to zero and not redistributed, so
that the shares of the zone model do not sum to one there; this was fixed in the protocol and its effect was
not assessed. The table holds the 2{,}548 pupils with an episode reported in the prospective survey of the
source (96\pct{} of the cases known to the schools) and the 8{,}375 who reported none in a later survey,
10{,}923 or 83\pct{} of the 13{,}217 enrolled; the missing pupils are mostly non-cases, so the denominators of
the classes are incomplete. The source excluded three small schools (71 pupils) from its own analysis; they
are kept here. The model has no household term, although the source attributes 41\pct{} of the pupils'
infections to households (Section~\ref{M-s8b_evidence:sec:zones}).
% src: code/bsc_validation2/a4_zone_level/PREREGISTRATION.md section 2.3; code/bsc_validation2/a4_zone_level/src/zonecore.py (features: Pc = Kc/(n-1), Pg = (Kg-Kc)/(ng-n), Ps = (Ks-Kg)/(ns-ng)); data/checks/zones.json: single_class_grades (36 grades, 453 pupils), respondents_share_of_enrolled_pct (82, as recorded; the 83 % of the text is 10,923/13,217 from the source, below) (notes/VERIFICATION.md, Section 4.9, zone level, further points 6, 7); Endo et al. 2021, Materials and Methods (Europe PMC PMC8609560, read 1 October 2026: 13,217 enrolled; 2,548 episodes, 96 % of the cases recognised by the schools; 8,375 without influenza in the retrospective survey; 10,923 combined; 71 pupils of 3 schools excluded; gamma kernel mean 1.7, SD 1.0, mean serial interval 2.2 d)
Power was computed before the freeze on synthetic epidemics on the real class structure (40 replicates per
truth, two strengths of the school signal). The verdict of the rule was earned in 100\pct{} and 80\pct{} of
replicates when the zone model is true, in none when the school is well mixed or the classes independent, in
7.5\pct{} and 0\pct{} under the generic model with a ratio drawn log-uniformly between 1 and 1{,}000, and in 2.5
and 5\pct{} under random shares. In the re-check the generic model turned out to be a more dangerous
competitor than these averages suggest: with its ratio set to the fitted value, about 37, it earns the verdict
in 6 of 16 replicates. The tolerance of 0.15 is a judgement, two to three times the sampling noise under a
true zone model; held-out differences depend slightly on a lower bound of $10^{-9}$ on the external hazard,
which is active in one fold (Table~\ref{s8b_evidence:tab:zone}).
% src: data/v2_numbers.json (zone.power); data/checks/zones.json: generic_ratio_model_at_fitted_ratio (ratio 37, 6/16); notes/VERIFICATION.md (Section 4.9, zone level, items 13, 7, 3)

\subsection{Register of tests}\label{appS_prot2:sec:register}

Table~\ref{appS_prot2:tab:register} lists every test of Section~\ref{M-s8_outbreaks:sec} of the main text, and the secondary and sensitivity
analyses of their protocols that are reported there and in Sections~\ref{appS_first:sec}--\ref{appS_prot2:sec},
with their status and their outcome after the re-checks. `Protocol' means that the test and its criterion were fixed before any model was run on
the data concerned; `listed' that the analysis is on the list of secondary or sensitivity analyses of a
protocol, which may never upgrade a primary outcome; `after' that it was added after results were known, by
the analysis or, where marked, by its re-check.

{\footnotesize
\setlength{\tabcolsep}{3pt}
\setlength{\LTcapwidth}{\textwidth}
\begin{longtable}{@{}>{\raggedright}p{4.6cm}>{\raggedright}p{1.5cm}>{\raggedright\arraybackslash}p{9.4cm}@{}}
\caption{Register of the tests of Section~\ref{M-s8_outbreaks:sec} of the main text and Sections~\ref{appS_first:sec}--\ref{appS_prot2:sec}.}
\label{appS_prot2:tab:register}\\
% src: default-parameter row: data/s8_default_level.json (expected_infections_bound 0.1021, 0.1570; beta_needed_per_day 16.34, 13.61 / 11.99)
% src: each row summarises Section 7 of the main text and Sections S7 to S9, where its numbers carry their own src comments; M1 row: data/v2_numbers.json (out2.pooled.M1.all: delta0 15.446, deltaC -8.055, p_mirror 1.5e-5; out2.pooled.M1.air: deltaC -4.634; out2.pooled.M1.decision); attainable totals: data/bsc_validation/02b_m1_exact.json, code/bsc_validation/POSTHOC_LOG.md (P4); zone rows: data/v2_numbers.json (zone.D_Z_minus_C; zone.check.hazard_bound."1e-09".D_N702); status: the four protocols and post-hoc logs; added rows: data/v2_numbers.json (out2.sec.train_kernel.rooms, out2.sec.loo_air, out2.sec.loo_room, out2.sec.seat, out2.sec.rho2.4, out2.pooled.M3.all, out2.pooled.M3.decision, out2.variants.summary, out2.rr_pooled."without P1", out2.check.bins."5 bins {0},{1},{2},{3-5},{6+}".air.deltaC -0.777, tr.power.pattern); data/s8_first_test_exact.json (primary.pooled.M2, loo_room)
\toprule
Test & Status & Outcome after the re-check \\
\midrule
\endfirsthead
\toprule
Test & Status & Outcome after the re-check \\
\midrule
\endhead
\bottomrule
\endfoot
\multicolumn{3}{@{}l}{\emph{First outbreak test (Section~\ref{M-s8_outbreaks:sec:holdout} of the main text and Sections~\ref{s8_outbreaks:sec:calibration}--\ref{s8_outbreaks:sec:outcomedetail})}} \\
Flight VN54, near/far split & protocol & pass ($\Mtwo$ $p=0.61$, $\Mzero$ $p=0.0008$); 20 passengers; a constant ratio of 2 to 3 also passes \\
Classroom, front/back split & protocol & not a real prediction: rejected at the calibrated value ($p=0.0009$); $p=0.068$ over the posterior \\
Bus, zone split & protocol & compatible ($p=0.38$), non-discriminating \\
Coach, rear/front split & protocol & fail ($p=0.007$) \\
Call centre, desk-scale clustering & protocol (quadrature after, check) & fail (0.007 to 0.013 by quadrature; 0.013 to 0.024 with the prior extended to $10^{6}\sqmh$; with the re-check's mode-sum propagator; its round-off acts only through the outbreak-size weights, and with zero weight at or below $0.1\sqmh$ the second value would be 0.015; the 40-draw estimate of the protocol had Monte Carlo error large enough to read as borderline) \\
Pooled test against $\Mzero$ & protocol & weak pass ($+2.83$, $p=2.5\times10^{-4}$, after the numerical correction of the classroom; $+3.87$ before), resting on the flight; constant ratios of 2 to 3 score higher (after, check) \\
Classroom with the round-off-safe propagator & after & $p=0.068$ instead of 0.11; outcome W4 unchanged \\
Totals attainable by $\Mone$ & protocol (rule P4 after) & mean number of cases below the observed one on bus, flight and classroom; observed total reached with probability 0.0002, 0.025 and 0.081, so unattainable under rule P4 on the bus only \\
Level checks and negative controls & protocol & met, with almost no power \\
Common coefficient, four vehicle events & after & rejected (approximate $p=0.023$, LR 9.52); for the two rooms not rejected (approximate $p=0.080$) \\
Default parameters of Sections~\ref{M-s2_model:sec}--\ref{M-s6_scenes:sec} against the bus and the choir & after & expected infections at most 0.10 and 0.16 against 23 and 52 (32 confirmed) cases; reached only for $\beta$ of about 16 and 14 (12)$\perday$ \\
\addlinespace
\multicolumn{3}{@{}l}{\emph{Second outbreak test (Section~\ref{M-s8_outbreaks:sec:second})}} \\
B1: $\Mtwo$ against $\Mzero$, five held-out events & protocol & pass ($+14.35$) \\
B2: $\Mtwo$ against the constant ratio, pooled & protocol & fail ($-9.15$), nine-tenths from the processing line, whose near zone was chosen by the source on the same data \\
B2 on the four flights & protocol & inconclusive ($-0.77$; power 0.82) \\
B3: adequacy of $\Mtwo$ & protocol & fail on 2 of 5 events; outcome V5 \\
$\Mone$ in closed form & protocol & V5, as $\Mtwo$; by the protocol a pass of the closed form, which overstates the reach of $\Mone$, does not bear on $\Mone$: against $\Mzero$ pass ($+15.45$); against the constant ratio pooled $-8.06$ (mirror $p=1.5\times10^{-5}$), on the flights $-4.63$; fails on 3 of 5 events \\
Train kernel applied to eight flights & listed & fail: four flights rejected \\
Exchange of calibration and hold-out & listed & weak pass on flights ($+3.58$, $p=0.004$); fail in the ward \\
Train kernel applied to the rooms & listed & rejected on the ward inpatients ($p=2\times10^{-5}$) and the processing line ($p=4\times10^{-4}$) \\
Leave one event out within each class & listed & flights not separated ($\Delta_{\mathrm{C}}=-0.06$; $p=0.046$ under the constant ratio, mirror $p=0.044$ under $\Mtwo$: improbable under either); rooms no better than well mixed, constant ratio better ($-14.0$) \\
Seat-level log score, three held-out flights & listed & constant ratio $-76.4$, $\Mtwo$ $-78.6$, well mixed $-83.5$ \\
Published ratio 2.4 in place of the calibrated one & listed & flights $+1.18$ ($p=0.028$); pooled $-4.44$ (mirror $p=0.002$; room kernel as first frozen) \\
Two-scale variant ($\Mtwo$ plus a well-mixed fraction) & protocol, exploratory & V3: $+17.5$ against well mixed, $-5.95$ against the constant ratio \\
Common coefficient within a class & listed & rejected (approximate $p$: rooms $1.6\times10^{-4}$, flights $0.034$) \\
Fifteen sensitivity variants & listed & V5 in 12, V3 in 3 (kernel as first frozen); on flights $\Delta_{\mathrm{C}}$ from $-3.42$ (constant ratio significantly better in three) to $+1.21$ ($\Mtwo$ significantly better in one) \\
Near/far risk ratios & listed & description: 3.98 (2.32--6.82; random effects); no gradient for influenza \\
Near/far ratio without the processing line & after & 3.30 (1.99--5.45) \\
Numerical correction of the room kernel & after & outcome unchanged \\
Other near zones and strata; ratio in the form of the protocol & after, check & $\Delta_{\mathrm{C}}$ on flights from $-0.78$ to $+7.8$ in these specifications; with the listed variants the sign is not determined \\
Exponential kernel & after, check & as good as $\Mtwo$ on flights, better on the processing line \\
\addlinespace
\multicolumn{3}{@{}l}{\emph{Contact records (Section~\ref{M-s8b_evidence:sec:contacts})}} \\
Lattice walk, structure and epidemic outcomes, six datasets & protocol & fail (A0) on all six \\
Walk against models without space & protocol & adds nothing; a static network of pair rates does better (after, check) \\
Modified walks on three confirmatory datasets & protocol (second freeze) & formally A0; uninformative, since the fit does not recover its own truth (after, check) \\
Error sizes of the modified walks & after & description only \\
Density scaling of contacts & listed & walk: exponent about 2; data below 1 in four of five settings \\
Power of the rule & protocol & adequate for the walk, weak near the tolerance \\
\addlinespace
\multicolumn{3}{@{}l}{\emph{Tracer measurements (Section~\ref{M-s8b_evidence:sec:tracer})}} \\
Coach & protocol & not rejected; non-discriminating \\
Call centre & protocol & not rejected; non-discriminating; wrong sign of clustering \\
Bus & protocol & inconclusive \\
Classroom & protocol & weak pass; constant ratios do as well \\
Flight VN54 & protocol & not reproduced (formal $p=0.14$; rejected in four of six listed variants) \\
Restaurant & protocol & not reproduced at the reported count \\
Trains & protocol & fail; every model fails \\
Pooled rule & protocol & S4 as computed, S5 on the defensible reading of the flight; not better than constant ratios \\
Probability of the pooled pattern under each truth & after, check & 0.23 under the tracer kernel, 0.84 and 1.00 under ratios 2 and 3 \\
Measured against fitted coefficients & listed & vehicle values 4.5 to about 50 times the first test's vehicle value; second-test aircraft value inside the cabin range, which does not make it a measurement of air mixing \\
Four further flights & after & exploratory, outcome-aware: weak \\
\addlinespace
\multicolumn{3}{@{}l}{\emph{Zone level (Section~\ref{M-s8b_evidence:sec:zones})}} \\
Tier 1: against well-mixed school and independent classes & protocol & pass \\
Tier 2a: equality of shares & protocol & fail \\
Tier 2b: tolerance; tier 3: dispersion & protocol & weak pass \\
Against a generic model with one fitted ratio & protocol (not decisional) & not beaten: zone model ahead by 18.8 ($-8.8$ to 54.0) out of sample, generic model better by 13.6 on the full data \\
Against guessed shares & after, check & guess 0.7/0.1/0.2 better by 19 (3 to 38); guesses chosen after the fitted shares were known; school-level share confounded by households \\
Frequency- against density-dependent transfer & listed & frequency-dependent fits better; not re-run by the check \\
Building by floor & protocol (low power, not blind) & inconclusive \\
Wings of one floor; early growth; hotspot maps & -- & not run: no independent measurement of mixing found \\
\end{longtable}
}

\subsection{Data, licences and what is redistributed}\label{appS_prot2:sec:data}

\emph{Outbreaks.} The seat maps and stratum counts of the eleven events were transcribed by eye from the
published figures and tables into the code and are released as tables, with the reference and the table or
figure of origin of every number (\texttt{data/outbreaks2/}). Copies of the source texts and figures are not
redistributed. \emph{Contact records.} The SocioPatterns files are distributed by their authors under
Creative Commons licences, some of which exclude commercial use; they are not redistributed here. A script
downloads them from the source and checks their SHA-256 against the values recorded when the protocol was
frozen. \emph{Tracer measurements.} The cabin measurements are published on figshare under CC~BY~4.0
\citep{kinahan2021aerosol}; the released tables are derived from those workbooks. The values for the coach and
the carriage were read or digitised from the supplementary figures of the sources and are released as tables
of numbers. \emph{Schools.} The anonymised pupil-level data of \citet{endo2021within} are published by their
authors under the MIT licence; the release contains a tabular copy, with that licence, and the script that
produced it from the original file.
% src: code/bsc_validation2/a2_contact_mobility/data/raw_SHA256.txt; data/tracer/SOURCES.md; code/bsc_validation2/a4_zone_level/PREREGISTRATION.md section 2.1 (MIT licence, commit 658abfe)
